\documentclass[11pt]{article}
\usepackage[margin=1.03in]{geometry}
\usepackage{amsmath,amssymb,amsthm}
\usepackage{setspace}\usepackage{textcomp}
\usepackage{booktabs}
\usepackage{graphicx}
\usepackage{array}\usepackage{multirow}
\usepackage[hidelinks]{hyperref}
\usepackage{natbib}
\setstretch{1.12}
\newtheorem{hypothesis}{H}
\newcommand{\CE}{\mathrm{CE}}
\newcommand{\VaR}{\mathrm{VaR}}

\title{\textbf{The TIPS--Treasury Basis as a Limits-to-Arbitrage Equilibrium Floor}\\[5pt] \large Theory, Data, Methodology, and First Results\\ \normalsize Second Dissertation Paper --- Comprehensive Working Document\\[3pt] \small Version 3.1 --- 23 July 2026 (v3 incorporated the consolidated external defect review; v3.1 adds the law-of-one-price set-up and generalises the international extension to two markets)}
\author{Alessio Ottaviani\\ EDHEC Business School, PhD in Finance\\[4pt] \small Supervisor: Prof.\ Riccardo Rebonato}
\date{}

\begin{document}
\maketitle

\begin{abstract}
\noindent This document consolidates, in one place, the theoretical framework, the data, the full empirical methodology, and every result obtained to date for the second dissertation paper. The paper models the post-compression TIPS--Treasury basis not as a dying anomaly but as an \emph{equilibrium floor}: the minimum rent at which a risk-averse, capital-constrained arbitrageur is indifferent between entering a trade whose payoff is certain at maturity and holding the capital that the trade's forced-unwind risk consumes. The theory section formalises the supervisor's framework --- mark-to-market dynamics, a Value-at-Risk unwinding barrier, a first-passage mixture for terminal wealth, a certainty-equivalent entry condition, and the comparative statics it delivers --- including a heterogeneous-threshold extension whose distinctive prediction is cross-sectional \emph{dispersion} in stress. The empirical sections report: a construction validated three independent ways and now certified pair-by-pair against Fleckenstein--Longstaff--Lustig's own Table~III (long pairs within $\pm3$\,bp); a post-2013 floor of 16--24\,bp with Newey--West $t$-statistics up to 24 and 46.8\,bp on the exact FLL window against their 54.5; a \emph{staircase} compression (breaks December 2010, sup-$F=160$, and December 2017) with a 2024 closure of the short-end basis; Ornstein--Uhlenbeck dynamics with a 3.7--4.6-month half-life; a 52\,bp CPI seasonal at two years; three-to-seven-fold conditional stress moves with multi-month lags and zero unconditional linear volatility loading --- the barrier signature; March 2020 identified as a cross-sectional dispersion event; an executed constraint test in which a volatility shock loads with coefficient 7.15 [$t=4.1$] when the He--Kelly--Manela capital ratio is in its bottom quintile and is small and at best marginal otherwise, with the episode-sorting table (GFC at the 3rd HKM percentile, COVID at the 8th, versus 32nd/37th for 2022/SVB) as the discriminating exhibit the formal threshold tests honestly underpowered (sup-$F$ $p\approx0.25$; pre-specified-quintile Chow $p=0.10$/$0.15$), and a leave-one-episode-out check disclosing that the interaction is substantially GFC-driven --- the disclosed reason the euro extension is decisive; a vintage-by-buffer forced-unwind surface; a corrected structural pass --- daily barrier monitoring aligned with the replay, the note's exact hold-to-maturity mixture restored, the stopped branch discounted, and the certainty-equivalent entry condition implemented --- under which, at the $\approx$2\% capital charge that the Basel $m\times\VaR$ rule itself \emph{derives} from the estimated dynamics (2.26\%/1.91\%), the observed floor of 17.7--23.3\,bp corresponds to relative risk aversion between one and two; and a notional-weighted replication of FLL's Figure~2 with a pair-level diagnostic that localises the residual measurement gap to the short-maturity pairs and identifies its two causes. Two appendices complete the record: a reconciliation with the June 2026 Detailed Project Description, and a register of every test executed to date with its script and headline output.
\end{abstract}

\tableofcontents

\section{Introduction and the Delimited Claim}

A Treasury Inflation-Protected Security combined with a zero-coupon inflation swap of matching maturity replicates, state by state, the payoff of a nominal Treasury. \citet{fll2014} documented that between 2004 and 2009 the synthetic traded persistently cheap --- an average mispricing of 54.5 basis points, above 200 at the crisis peak --- and attributed its time variation to slow-moving arbitrage capital. Since then the basis has compressed to a narrow band, and the literature's default reading is that of a closing anomaly. This paper defends a different reading: the basis has converged not to zero but to a small, stable, strictly positive \emph{resting level} that widens violently --- and briefly --- in episodes of intermediary distress, and this resting level is an economic object in its own right: the equilibrium rent of the marginal, capital-constrained arbitrageur.

The delimited claim, defended against each adjacent literature in Section~\ref{sec:lit}, is that \emph{no existing paper derives the equilibrium resting level of the TIPS--Treasury basis from a structural model with endogenous forced unwinding and a capital-charge participation constraint, and confronts its level, term-structure, stress, dispersion, and survival-surface predictions with the data}. The empirical work summarised here shows the claim is not only unoccupied but, on first confrontation, quantitatively promising: with dynamics estimated from the data and the capital charge derived from the regulatory rule rather than assumed, the observed floor corresponds to relative risk aversion between one and two.

\section{Theory}\label{sec:theory}

\subsection{The no-arbitrage identity and the basis}

Following the supervisor's note, denote nominal quantities in upper case ($P_0^T$, $Y_0^T$) and real quantities in lower case ($p_0^T$, $y_0^T$), with $f_0^T$ the breakeven inflation rate and $\pi(t)$ instantaneous inflation. A unit of TIPS pays $\exp[\int_0^T\pi(s)ds]$ at maturity and the inflation swap pays $\exp[f_0^T\,T]-\exp[\int_0^T\pi(s)ds]$, so the portfolio $\Pi_0^T$ of one TIPS plus one swap pays $\exp[f_0^T\,T]$ with certainty in every state of the world. Since entering the swap is costless, the law of one price requires that $\exp[-f_0^T T]$ units of the combination cost exactly the nominal zero-coupon bond:
\begin{equation}
\exp\!\left[-f_0^T\,T\right]\cdot p_0^T \;=\; P_0^T
\qquad\Longrightarrow\qquad
p_0^T \;=\; P_0^T\cdot\exp\!\left[+f_0^T\,T\right],
\label{eq:loop}
\end{equation}
consistently with $P_0^T\equiv\exp[-(y_0^T+f_0^T)T]$ and $p_0^T=\exp[-y_0^T\,T]$.\footnote{The rearrangement carries the \emph{positive} exponent. Equation~(4) of the supervisor's note prints $p_0^T=P_0^T\exp[-f_0^T T]$, which contradicts both its own left-hand side and the note's eqs.~(5)--(6), whose ratio is $\exp[+f_0^T T]$; economically, the real bond carries the lower yield and therefore the higher price. The slip is a dead end within the note --- nothing downstream depends on it, and eqs.~(5)--(6) restate the correct relation independently --- and it has no bearing on any result in this paper. It is recorded here alongside the eq.~(16) correction of Section~2.2.} When the portfolio instead trades \emph{below} the nominal bond it does so at a higher yield, $Y(\Pi_0^T)=Y_0^T+\lambda_0^T$ with $\lambda_0^T>0$: the object of the paper. For the constant-maturity implementation used throughout, let $y^{\text{nom}}_t(\tau)$, $y^{\text{real}}_t(\tau)$ and $s_t(\tau)$ denote the nominal yield, the TIPS real yield, and the zero-coupon inflation-swap fixed rate at tenor $\tau$; holding the TIPS and paying fixed on the swap converts the real payoff into a certain nominal one, the synthetic nominal yield is $y^{\text{real}}_t+s_t$, and the basis
\begin{equation}
\lambda_t(\tau) \;=\; \big(y^{\text{real}}_t(\tau)+s_t(\tau)\big)-y^{\text{nom}}_t(\tau) \;=\; s_t(\tau)-b_t(\tau),
\label{eq:basis}
\end{equation}
with $b_t$ the breakeven rate, is the yield pick-up of the synthetic over the actual nominal. Under the law of one price $\lambda=0$; the observed $\lambda>0$ is the object of the paper. Equation~\eqref{eq:basis} is also the operational construction: it requires only three curve points per maturity, is computable daily, and --- as Section~\ref{sec:valid} documents --- tracks the exact bond-level construction closely at the benchmark maturities.

\subsection{The strategy's mark-to-market}

Following the supervisor's note, the arbitrage buys the cheap synthetic, sells the nominal, and invests the net cash at the nominal yield. Every leg is a transparent market price, so the strategy has zero value at inception; with $\tau\equiv T-t$ its time-$t$ value is
\begin{equation}
S(t) \;=\; e^{-Y_0^T\tau}\big(1-e^{-\lambda_0^T T}\big)\;-\;e^{-Y_t^T\tau}\big(1-e^{-\lambda_t^T\tau}\big),\footnote{The first term carries the \emph{full-horizon} spread exponent $\lambda_0^T T$: with $\lambda_0^T\tau$, as in eq.~(16) of the supervisor's note (whose eqs.~11, 12 and 14 carry the correct form), the expression would give $S(T)=0$, contradicting the certain terminal payoff. The correction restores $S(0)=0$ and $S(T)=1-e^{-\lambda_0^T T}$; it has no empirical impact here (all tables use the linear approximation) and is to be confirmed with the supervisor.}
\label{eq:mtm}
\end{equation}
the difference between a deterministic accrual and a stochastic term in the nominal yield $Y_t^T$ and the spread $\lambda_t^T$. At maturity $S(T)=1-e^{-\lambda_0^T T}>0$ with certainty: the trade is an arbitrage \emph{if it can be held}. Before maturity a widening spread or falling yields push $S(t)$ negative --- and the two move adversely together in flight-to-quality states. For empirical work the first-order approximation $S(t)\approx -D\,(\lambda_t-\lambda_0)$ in basis points of notional, with $D$ the trade's duration, is used throughout and is the object the forced-unwind replay tracks.

\subsection{Dynamics, the barrier, and the terminal mixture}

The spread follows a mean-reverting Ornstein--Uhlenbeck process under the physical measure,
\begin{equation}
d\lambda_t \;=\; \kappa\,(\theta_t-\lambda_t)\,dt+\sigma_t\,dW_t,
\label{eq:ou}
\end{equation}
consistent with the strong mean reversion the data display (Section~\ref{sec:dyn}), with the yield factor correlated Gaussian in the full model. The arbitrageur faces a risk-management constraint: the position is unwound the first time the mark-to-market loss exceeds a multiple of its capital charge,
\begin{equation}
\tau_B \;=\; \inf\{t:\ -S(t)\ \ge\ m\cdot\VaR_\alpha\}\;\approx\;\inf\{t:\ D(\lambda_t-\lambda_0)\ \ge\ B\},
\end{equation}
with $B$ the capital buffer in P\&L terms. Terminal wealth is therefore a \emph{mixture}: a mass point at the certain carry, with weight equal to the first-passage survival probability $\Pr(\tau_B>T)$, plus a continuous distribution of stopped outcomes $\{\text{carry to }\tau_B - B\}$. For the OU specification survival probabilities are available semi-analytically \citep{alili2005,linetsky2004}; the Monte Carlo implementation is the workhorse and, in Section~\ref{sec:struct}, the estimation engine.

\subsection{The entry condition and the equilibrium floor}

An arbitrageur with utility $u$ (exponential in the baseline) enters iff the certainty equivalent of terminal wealth exceeds committed wealth, identified with the capital charge:
\begin{equation}
\CE\big[\,W_0+X_T(\lambda_0;\kappa,\theta,\sigma,D,B,T)\,\big]\ \ge\ W_0,\qquad W_0=\text{capital charge}.
\label{eq:entry}
\end{equation}
Free entry drives the basis down until \eqref{eq:entry} binds, defining the floor $\lambda^*=\inf\{\lambda_0:\CE\ge W_0\}$: below $\lambda^*$ no capital enters and the basis cannot compress further. Two implementation points, both surfaced by the external defect review and both now implemented, matter to first order. \emph{Horizon.} The note's terminal mixture --- a mass point at the certain carry plus a stopped branch --- exists only when the position is held to bond maturity $T$; at a shorter risk horizon $H$ the surviving branch carries the mark-to-market law of $\lambda_H$, not a certainty, and at $H=24$ months with a ten-year trade that law has a standard deviation of the same order as several years of carry. Section~\ref{sec:struct} therefore runs both variants: the maturity engine, which restores the note's mixture exactly, and the $H$-horizon engine with the MTM law for survivors. \emph{Risk aversion.} Because the surviving branch is riskless only at maturity, the CARA adjustment is first-order in the $H$-horizon variant --- the mean--variance penalty $\tfrac{\gamma}{2}\mathrm{Var}[X]$ with $\gamma=\text{RRA}/W_0$ consumes 34--137\% of $\mathbb{E}[X]$ at the estimated terminal law --- so the certainty equivalent is computed exactly on the simulated terminal distribution and the floor is reported as a function of risk aversion, with the risk-neutral participation floor and a 10\%-hurdle version retained as reference points. The normalisation $\gamma=\text{RRA}/W_0$ (curvature over committed capital) is a stated choice; curvature over total firm wealth would shrink the penalty, and the floor-versus-RRA curve makes the mapping transparent either way. Finally, the capital charge $W_0$ itself is not asserted: Section~\ref{sec:struct} \emph{derives} it from the note's $m\times\VaR_\alpha$ under the estimated dynamics.

\subsection{Heterogeneous thresholds and the dispersion prediction}

Real-world arbitrage capital is not a representative agent: desks differ in buffers, funding terms and mandates. Let the buffer be distributed as $B\sim F$ across arbitrageurs, with holdings matched assortatively to bonds by balance-sheet cost. The aggregate unwind mass after a shock of size $s$ is then
\begin{equation}
\Pi^{\text{unwind}}(s) \;=\; \int \Pr\!\big(\tau_{B(b)}\le T \mid s\big)\, dF(b),
\label{eq:mixture}
\end{equation}
and the model's distinctive cross-sectional prediction follows: a common shock stops the most constrained holders of the most balance-sheet-expensive bonds \emph{first}, so the cross-section of bond-level bases \emph{fans out before the median moves}. Dispersion --- the interquartile range of the CUSIP-level basis --- is therefore a first-class model object, its stress behaviour a signature unavailable to representative-agent alternatives, and Section~\ref{sec:disp} shows March 2020 delivering it on cue. The mixture also supplies the policy-counterfactual apparatus carried over from the June project description: shifting $F$ (a leverage-ratio recalibration; a change in haircut regimes) moves both the floor and the dispersion response in computable ways, and the counterfactual section of the final paper will report the implied floor under alternative regulatory calibrations of $(m,\alpha)$ and $F$.

\subsection{Comparative statics and testable predictions}

\begin{hypothesis}\textbf{Existence.} Post-compression, $\lambda$ is stationary and mean-reverting around a strictly positive $\theta$: the floor exists.\end{hypothesis}
\begin{hypothesis}\textbf{Term structure.} $\lambda^*(\tau)$ rises with maturity over the segment where path risk per unit of carry rises; maturity-dependent frictions (haircuts, term-repo scarcity) can flatten or locally invert it.\end{hypothesis}
\begin{hypothesis}\textbf{State dependence.} The floor loads on the \emph{constraint}: linearising \eqref{eq:entry}, $\theta_t=\theta_0+\xi H_t$ with $H_t$ intermediary-capital scarcity and $\xi>0$; shocks move the basis \emph{when the constraint binds} and barely otherwise --- large conditional responses with a near-zero unconditional linear loading is the model's signature, not its failure.\end{hypothesis}
\begin{hypothesis}\textbf{Regime shifts.} Discrete changes in the capital regime (Basel phase-ins, leverage-ratio recalibrations, the growth of arbitrage capital) shift the resting level in steps.\end{hypothesis}
\begin{hypothesis}\textbf{Survival surface.} The model's first-passage probability maps one-for-one into the empirical stop-out frequency by entry vintage and buffer --- a dense, over-identifying set of calibration moments.\end{hypothesis}
\begin{hypothesis}\textbf{Dispersion.} Per \eqref{eq:mixture}: the cross-section fans out in stress; dispersion is priced by the heterogeneity of $F$, not by the representative buffer.\end{hypothesis}

\subsection{Implementation choices and their signed effects}

Five further choices, each flagged in the defect review, are recorded with the sign of their effect. \emph{Marginal versus standalone VaR} (raised in the note's closing line): the engine charges standalone VaR; in a multi-strategy book the marginal charge is smaller, so the true floor is \emph{lower} --- standalone flatters the fit from above, and a stylised two-strategy appendix will bound the wedge. \emph{The omitted yield leg}: the note's $S(t)$ has two stochastic terms with adverse comovement in flight-to-quality; the current engine works in $\lambda$ alone. The review's recomputation on the exact $S(t)$ ($\sigma_Y=23.9$\,bp/m, $\rho=-0.135$ post-2013, $-0.315$ in 2008--09) puts the omission at $+3.5$pp on the 1\%-buffer stop-out ($+5.6$pp at crisis correlation) --- adopted provisionally here, to be verified in the bivariate engine of the SMM merge. \emph{Discounting}: the stopped branch is now discounted at $Y_0$ (a 3--4\% effect, previously omitted). \emph{Linearisation}: empirical objects use $S\approx-D\Delta\lambda$; a robustness line evaluating the exact eq.~\eqref{eq:mtm} at the 150--320\,bp GFC moves is scheduled. \emph{Zero-coupon theory versus coupon CUSIPs}: the panel's bonds carry coupons, duration drift, and floor moneyness heterogeneity that the zero-coupon algebra abstracts from; the cross-sectional median damps but does not eliminate these wedges, and the bond-level engine of the FLL replication is where they are resolved.

As $\sigma\to0$ the floor collapses to a margin/funding-cost bound of the \citet{gp2011} type; the volatility and state-dependence loadings of H3 are the data's handle for discriminating the path-risk mechanism from the pure funding-cost alternative.

\subsection{Relation to the literature}\label{sec:lit}

\citet{fll2014} document the mispricing and its slow-moving-capital variation; \citet{dkw2018} and \citet{cg2022} price TIPS liquidity; \citet{hns2022} dissect March 2020 --- none models the resting level. \citet{gp2011} and \citet{tuckmanvila1992} bound mispricings by funding costs without path risk; \citet{shleifervishny1997}, \citet{liulongstaff2004}, \citet{kondor2009} and \citet{gromb2010} supply the forced-unwind mechanism without confronting a specific basis quantitatively; \citet{hkm2017} and \citet{duffie2010} supply the state variables; \citet{mp2012} the arbitrage-crash phenomenology; \citet{fl2024} document the opposite-signed ``Treasury richness'' episodes that motivate the signed extension held in reserve. The contribution sits at the intersection: the structural floor, estimated and tested on the cleanest deterministic-payoff arbitrage in existence.

\section{Data}\label{sec:data}

\begin{table}[h!]
\centering\small
\begin{tabular}{p{4.1cm}p{6.6cm}p{3.6cm}}
\toprule
Block & Content & Coverage \\
\midrule
Basis inputs (Bloomberg) & USSWIT1--30; USGGT02/05/10/20/30Y; exact nominals USGG2/5/10/30YR + 20YR & daily; swaps 2004-07+, TIPS 1998+ \\
Cross-checks & USGGBE02--30 breakevens; 107-CUSIP bond-level basis panel with per-issue Amount Issued (notional weights) & daily; 2004-08 -- 2026-05 \\
FLL benchmark & Table III of \citet{fll2014}: the 29 pairs with per-pair mean/SD/min/max/$N$ (extracted; \texttt{fll\_pairs\_tableIII.csv}) & Jul-2004 -- Nov-2009 \\
Stress proxies & MOVE, VIX; LIBOR--OIS, SOFR/BGCR/TGCR; 8 dealer CDS; CDX/iTraxx & daily; 1998/2001+ \\
Constraint variables & HKM capital ratio (monthly, vintage 2025-06); FR\,2004 dealer TIPS net positions (weekly, spliced over six reporting breaks); CFTC TFF leveraged-funds Treasury DV01; OFR Form-PF repo and RV-NAV (quarterly) & 1970+/1998+/2013+ \\
Seasonality & CPURNSA (CPI-U NSA) & monthly, 1990+ \\
Public curves & GSW nominal and TIPS (Fed Board) & daily, 1961/1999+ \\
\bottomrule
\end{tabular}
\caption{Data inventory. Not available anywhere, and flagged as such: security-level TIPS repo specialness and bond-level bid--ask, which gate the transaction-cost band and the identification of the 2021 negative episode.}
\end{table}

\section{Basis Construction and Triple Validation}\label{sec:valid}

The basis is computed from \eqref{eq:basis} at $\tau\in\{2,5,10,20,30\}$ with \emph{exact} constant-maturity legs (2/5/10/30 from 2004-07; 20 from May 2020, the sector's reintroduction). Validation is threefold, and the full breakeven cross-check is reported by maturity:

\begin{table}[h!]
\centering\small
\begin{tabular}{lcccc}
\toprule
Maturity & Mean diff, direct vs BE (bp) & SD (bp) & Correlation & $n$ (daily) \\
\midrule
2Y  & $-3.97$ & 9.78 & 0.974 & 5648 \\
5Y  & $-5.88$ & 9.76 & 0.839 & 5718 \\
10Y & $\mathbf{-0.55}$ & \textbf{1.55} & \textbf{0.996} & 5718 \\
20Y & $-1.80$ & 6.50 & 0.686 & 1587 \\
30Y & $+0.60$ & 4.70 & 0.969 & 5718 \\
\bottomrule
\end{tabular}
\caption{Construction cross-check: direct \eqref{eq:basis} versus the Bloomberg-breakeven route $\lambda^{BE}=(s-\text{BE})\times100$. Certified at the benchmark maturities; the 4--6\,bp short-end divergence is the first quantitative signal that the short end requires bond-level treatment (Section~\ref{sec:fllrep}).}
\end{table}

Second, against the 107-CUSIP bond-level median: daily level correlation \textbf{0.926}, monthly-change correlation 0.591, regime means aligned within the $\sim$5\,bp curve-versus-bond wedge, GFC peaks 153.4 versus 150.6\,bp (Figure~\ref{fig:cusip}). Third, an independently run parallel reconstruction on the same legs matches to rounding. Measurement risk is closed: the floor and the blow-outs are facts about the market, and Section~\ref{sec:fllrep} pushes the certification to the individual-pair level against FLL's own published statistics.

\begin{figure}[h!]\centering
\includegraphics[width=0.9\textwidth]{fig9_cusip_vs_curve.png}
\caption{Bond-level median (107 CUSIPs) against the exact-leg curve construction: correlation 0.926, common regimes, common peaks.}\label{fig:cusip}
\end{figure}

\section{Empirical Methodology}\label{sec:method}

\textbf{(i) Level dynamics.} ADF tests and the discrete OU $\lambda_t=a+b\lambda_{t-1}+\varepsilon_t$, whence $\kappa=-\ln b/\Delta t$, $\theta=a/(1-b)$, half-life $\ln2/\kappa$; block-bootstrap bands. \textbf{(ii) Breaks.} Single mean-shift by SSR grid (10\% trimming) with sup-$F$, second break on the longer segment; the staircase is the object, not a maintained single step. \textbf{(iii) Seasonality.} Month-of-year dummies (post-2010), joint HAC $F$, amplitude, ADF on the deseasonalised series; joint seasonal-and-break estimation flagged where level shifts contaminate the dummy fit. \textbf{(iv) Events.} Pre-level $\to$ peak $\to$ half-reversion, per maturity, with the swap-leg dislocation caveat for March 2020 at ten years. \textbf{(v) Dispersion.} Cross-sectional IQR and upper percentiles of the 107-CUSIP panel per episode. \textbf{(vi) The frozen constraint design.} $y_t=\text{med}_{t+3}-\text{med}_{t-1}$ on shock $\Delta S_t$, split by the \emph{lagged constraint state}; $\theta$ by grid; sup-$F$ with wild and block-wild bootstrap; contemporaneous designs are dead on arrival (the documented multi-month lag) and are not run. \textbf{(vii) Forced-unwind replay.} Monthly entry vintages, 24-month horizon, $D=8$, stop at $D(\lambda_t-\lambda_0)>B$; stop-out frequency by vintage $\times$ buffer. \textbf{(viii) Structural estimation.} Stage one: \eqref{eq:ou} by ML per regime. Stage two: friction parameters by minimum distance on the survival surface (vii) \emph{plus} the calm/stress loading pair from (vi), with block-bootstrap moment covariance in the \citet{gkr2014} spirit for correctly sized over-identification tests; the CE refinement layered on the hurdle version, and the full SMM engine inherited from the June codebase (\texttt{tips\_treasury\_arb/}) as the implementation vehicle. \textbf{(ix) FLL replication.} Notional-weighted cross-sectional average per their Figure~2, with Table~III as a 29-row per-pair validation suite.

\section{First Results}\label{sec:results}

\subsection{The floor, and the FLL benchmark}

\begin{table}[h!]
\centering\small
\begin{tabular}{lccccc}
\toprule
Period & 2Y & 5Y & 10Y & 20Y & 30Y \\
\midrule
FLL window 2004/07--2009/11 & 44.0\,[2.9] & 34.6\,[9.2] & \textbf{46.8\,[7.5]} & --- & 51.0\,[8.5] \\
2010--2012 & 22.1\,[6.0] & 15.4\,[5.3] & 33.7\,[15.0] & --- & 37.9\,[31.1] \\
2013--2019 & 30.1\,[9.5] & 17.9\,[12.8] & 26.3\,[24.2] & --- & 29.6\,[12.0] \\
2020--2026 & 15.9\,[3.6] & 13.3\,[5.5] & 20.3\,[15.6] & 12.7\,[8.9] & 17.9\,[21.6] \\
Post-2013 (all) & 23.3\,[7.6] & 15.7\,[10.9] & \textbf{23.4\,[22.4]} & 12.7\,[8.9] & 24.0\,[13.3] \\
\bottomrule
\end{tabular}
\caption{Monthly means of $\lambda$ (bp), Newey--West $t$ in brackets. The ten-year proxy averages 46.8\,bp on the exact FLL sample against their bond-level 54.5. Fifteen years and the largest capital inflow in the trade's history later, the basis rests at 16--24\,bp, not zero.}
\end{table}

The term structure on exact legs reads 23.3 / 15.7 / 23.4 / 24.0 at 2/5/10/30 post-2013: a rising and highly significant 5$\to$10 segment, a flat top, and two quantified end-artefacts --- the 2-year elevation is the CPI seasonal (52\,bp amplitude, Section~\ref{sec:dyn}); the 20-year depression is the cheap-sector effect of the 2020 reintroduction. H2 reads as supported on the clean segment, with the friction-profile refinement live.

\begin{figure}[h!]\centering
\includegraphics[width=0.94\textwidth]{fig5_basis_exact.png}
\caption{The basis on exact legs, 2004--2026: the FLL plateau, the GFC spike, the staircase, the 2024--25 short-end closure.}
\end{figure}

\subsection{The staircase}

\begin{table}[h!]
\centering\small
\begin{tabular}{lcccc}
\toprule
Maturity & Break 1 & Means (bp) & Break 2 & Means (bp) \\
\midrule
2Y & 2009-10 & $44.4\to23.2$ (16.4) & 2024-10 & $25.9\to0.3$ (27.0) \\
5Y & 2009-11 & $34.8\to15.6$ (108.3) & 2024-10 & $16.7\to6.1$ (20.0) \\
10Y & \textbf{2010-12} & $\mathbf{45.8\to24.4}$ (\textbf{159.6}) & \textbf{2017-12} & $\mathbf{29.0\to20.7}$ (132.6) \\
30Y & 2014-12 & $46.0\to20.7$ (237.7) & 2019-09 & $24.7\to18.0$ (117.0) \\
\bottomrule
\end{tabular}
\caption{Break adjudication (sup-$F$ in parentheses). The dominant ten-year break is December 2010, not 2013; the compression is a staircase $46\to29\to21$\,bp. New fact: after October 2024 the two-year basis mean is statistically zero --- the short-end arbitrage has closed.}
\end{table}

\begin{figure}[h!]\centering
\includegraphics[width=0.9\textwidth]{fig7_breaks10y.png}
\caption{The ten-year basis with estimated regime means: the staircase, H4's raw material.}
\end{figure}

\subsection{Dynamics and seasonality}\label{sec:dyn}

Post-2013, the ten-year basis is stationary (ADF $p=0.008$) with an OU half-life of 4.6 months on the curve construction and 3.7 on the bond-level median ($\kappa=0.187$/m, $\theta=16.6$\,bp, $\sigma=5.8$\,bp/m); the two- and five-year proxies fail the ADF \emph{until} the month-dummy seasonal is removed --- month dummies fitted on the post-2010 monthly sample, jointly significant everywhere, amplitude \textbf{52.2\,bp} at two years (against a 23\,bp mean), 22.9 at five, 5.5 at ten --- after which the two-year monthly ADF flips to $p=0.000$ (post-2013 window). The 30-year non-stationarity is the 2014/2019 breaks sitting inside the window; within-regime it mean-reverts. One documented anomaly: dummy deseasonalisation over a window containing level shifts absorbs trend at ten years; the joint seasonal-and-break estimator is the designated fix, and the raw ten-year result stands meanwhile.

\subsection{Stress: events, lags, and the linear emptiness}

\begin{table}[h!]
\centering\small
\begin{tabular}{llccc}
\toprule
Episode & Series & Pre $\to$ Peak (bp) & Peak date & Half-reversion \\
\midrule
GFC & 5Y / 10Y / 30Y & $\to$ 144 / 151 / 154 & 2008-11 -- 2009-01 & --- \\
March 2020 & 2Y & $-1.3\to68.2$ & \textbf{2020-03-19} & 2020-03-27 \\
March 2020 & 10Y & $19.4\to30.6$ & 2020-03-06 & (swap leg dislocated) \\
2022 hiking & 2Y / 5Y / 10Y & $\to$ 67.9 / 37.0 / 32.1 & 2022-06/09 & weeks \\
2023 (SVB) & 2Y / 10Y / 30Y & $\to$ 60.4 / 33.7 / 25.2 & \textbf{2023-07/08} & weeks \\
\bottomrule
\end{tabular}
\caption{Stress events on exact legs. Three-to-seven-fold moves in every episode; peaks lag the initiating shock by weeks to months (SVB: shock in March, peaks in July--August) --- the slow-moving-capital lag. At ten years the March-2020 spike is muted because the swap leg dislocated with the TIPS leg: that episode requires the bond-level construction at benchmark maturities.}
\end{table}

Against this, the \emph{unconditional linear} volatility loadings are empty: post-2013 local projections of the ten-year basis on MOVE shocks are insignificant at every horizon ($|t|\le1.2$), and the rolling resting level regressed on rolling MOVE gives $\beta=-0.014$ [$t=-0.4$]. Large conditional moves with zero linear loading is H3's signature: a linear regression averages a flat normal-times relation with a violent constraint-binding one and recovers nothing. The design consequence --- state-conditional specifications only --- is implemented in Section~\ref{sec:constraint}.

\subsection{March 2020 as a dispersion event}\label{sec:disp}

The heterogeneous-threshold mixture \eqref{eq:mixture} predicts that a funding shock is read first in the \emph{cross-section}: the most constrained holders of the most balance-sheet-expensive bonds are stopped out before the centre of the distribution moves. March 2020 delivers this signature with unusual cleanliness, and the episode table below reports it under a uniform, stated convention for the pre-event level --- the 60-trading-day trailing mean of the cross-sectional IQR, ending five business days before the episode onset (onsets: 2008-09-15, 2020-02-20, 2022-03-16, 2023-03-09); the earlier single-day snapshots are retired. The restated numbers sharpen rather than weaken the pattern: the corrected 2022 pre-level (16.2 against the stale 23.3) makes the 2022 dispersion move \emph{larger}, and the SVB episode is confirmed as a non-event on this margin. The precise March-2020 statement is proportional and temporal, not absolute: from 14 February to the 16 March apex the median \emph{doubles} ($17.8\to34.6$\,bp, $+17.4$) while the IQR \emph{triples} ($5.3\to18.4$), and the timing separates them --- the IQR peaks on the dash-for-cash apex itself, the median's maximum (36.2) arrives only in late June. Finally, the review's maturity-confound objection is met with the within-bucket decomposition: \emph{inside} the under-five-year bucket the IQR explodes from $\sim$5.3 to \textbf{43.8}\,bp (peak 19 March), against $1.3\to4.4$ within the over-five-year bucket --- the fanning is holder- and security-level heterogeneity \emph{within} maturity class, not a maturity-axis artefact, with assortative matching of constrained holders to expensive bonds stated as the identifying assumption.

\begin{table}[h!]
\centering\small
\begin{tabular}{lccc}
\toprule
Episode & IQR pre $\to$ peak (bp) & Peak date & Median \\
\midrule
GFC & $19.0\to117.2$ & 2008-11-24 & blows out (153) \\
March 2020 & $5.6\to18.4$ & \textbf{2020-03-16} & doubles (34.6 on 03-16; max 36.2 late June) \\
2022 & $16.2\to35.2$ & 2022-08-26 & moderate \\
2023 (SVB) & $10.1\to14.1$ & 2023-03-06 & no dispersion event \\
\bottomrule
\end{tabular}
\caption{Cross-sectional dispersion (107 CUSIPs); pre-event levels under the uniform 60-day trailing rule stated in the text. The 90th percentile of the cross-section reaches 73.3\,bp in March 2020. Maturity split confirms short-end concentration (COVID maxima: 49.1\,bp under five years vs 31.6 over; GFC: 210 vs 162). Monthly $\Delta$IQR regressions on MOVE, linear and threshold, are weak ($|t|\le0.8$): dispersion, like the level, is event-concentrated.}
\end{table}

\begin{figure}[h!]\centering
\includegraphics[width=0.86\textwidth]{fig10_mar2020_bondlevel.png}
\caption{March 2020 at bond level: the median, the interquartile band, and the swap proxy.}
\end{figure}

\subsection{The constraint test}\label{sec:constraint}

The design was frozen on 15 July 2026 (Pilot Report v3, \S4.3) --- lag structure, threshold grid, bootstrap scheme --- \emph{before} any constraint series had been inspected; a full specification curve over lag, grid and shock definitions is scheduled as robustness. \emph{The market-proxy passes (recorded for completeness).} Two preliminary passes preceded the constraint test proper, both on the frozen lag-aware dependent variable. The \emph{contemporaneous} design ($\Delta\text{med}_t$ on $\Delta S_t$) found nothing on any market proxy (all bootstrap $p>0.14$, perverse signs) and is retired. The \emph{lag-aware} design on market-price proxies produced: dealer CDS $b_{\text{calm}}=0.488$ [3.0] / $b_{\text{stress}}=-0.122$ [$-0.9$], sup-$F$ 38.2, $p=0.023$ but with the \emph{perverse} pattern (loading dies above the threshold, the 2008--09/2011--12 European-names artefact); LIBOR--OIS 0.143 [3.3] / 0.517 [3.4], $p=0.21$; MOVE 0.055 [1.2] / 0.549 [2.9], $p=0.36$. Lesson: market-price proxies are \emph{shocks}, not constraints --- they peak and revert before the basis does --- which motivated the move to the constraint variables below.

\begin{table}[h!]
\centering\small
\begin{tabular}{llccccc}
\toprule
Spec & Shock $\mid$ state (stress side) & $\hat\theta$ & $b_{\text{calm}}$ [$t$] & $b_{\text{stress}}$ [$t$] & sup-$F$ & boot $p$ \\
\midrule
A & $\Delta$HKM $\mid$ HKM low & 4.57\% & $-3.91$ [$-2.5$] & $\mathbf{-24.51}$ [$-2.6$] & 6.8 & 0.24--0.26 \\
B & $\Delta$MOVE/10 $\mid$ HKM low & 4.81\% & 0.63 [1.7] & $\mathbf{7.15}$ \textbf{[4.1]} & 21.8 & 0.25--0.26 \\
C & $\Delta$PD $\mid$ PD-z high & 1.90 & $-0.30$ [$-0.8$] & 1.02 [1.4] & 0.9 & 0.67 \\
C2 & $\Delta$MOVE/10 $\mid$ PD-z high & 1.16 & 4.93 [2.5] & $-1.55$ [$-0.8$] & 17.2 & 0.10 \\
D & $\Delta$TFF $\mid$ TFF-z low & $-1.46$ & $-0.08$ [$-2.0$] & 0.01 [0.3] & 1.9 & 0.23 \\
\bottomrule
\end{tabular}
\caption{The frozen lag-aware design on the constraint variables (2004-08--2026, monthly; bootstrap $p$ ranges across iid- and block-wild schemes, which agree). A volatility shock moves the basis eleven times harder when intermediary capital is in its bottom quintile; capital shocks six times harder. Inventories (C) and futures positioning (D) carry no threshold information: the binding constraint is capital. The sup-$F$ does not formally reject linearity ($p\approx0.25$) --- a power constraint with $\sim$50 stress months, stated as a referee would. Quarterly Form-PF interactions (hedge-fund repo $\times$ RV-NAV state) are uninformative at $n=52$, as anticipated. Three qualifications from the defect review are adopted. The ``eleven-fold/six-fold'' ratios are descriptive, not inferential --- their denominators are the calm loadings, small and at best marginally significant --- so the claim is stated as \emph{a significant stress-state loading against a calm loading near zero}. The near-significant perverse Spec C2 ($p=0.10$, sign flip) has a ready reading: high dealer TIPS inventories mark QE-absorption states, not stress, so the conditioning variable is inverted --- an interpretation, flagged as such. And because the objects are event-concentrated, permutation/placebo-window inference is scheduled alongside the bootstrap.}
\end{table}

Because the sup-$F$ prices the search over the threshold, and the bottom quintile was \emph{pre-specified} by the model (H3), the review's suggestion of a Chow test at the known quintile is implemented: Spec A returns $F=5.5$, block-wild $p=0.097$; Spec B returns $F=21.8$, $p=0.145$ --- an improvement on the sup-$F$'s 0.25, marginal for A, still short of conventional levels for B. Two disclosures complete the honest record. The stress-state \emph{composition}: the bottom-quintile months are 2008--2013 (42 months: the GFC aftermath and the euro crisis) plus 2020 (8 months) --- two macro impairment clusters in twenty-two years. And the \emph{leave-one-episode-out} check: excluding the COVID year leaves the stress loading intact (7.28 [$t=4.1$]); excluding the GFC window collapses it (2.46 [$t=1.8$]). The interaction is identified substantially off the GFC cluster --- which is precisely why the euro cross-section, supplying a third impairment episode, is the decisive addition rather than a nicety.

\begin{table}[h!]
\centering\small
\begin{tabular}{lccc}
\toprule
Episode onset & HKM level & Percentile (2004+) & Bottom quintile \\
\midrule
GFC (Oct-2008) & 3.57\% & \textbf{3\%} & yes \\
COVID (Mar-2020) & 4.27\% & \textbf{8\%} & yes \\
2022 (Jun-2022) & 5.41\% & 32\% & no \\
SVB (Mar-2023) & 5.62\% & 37\% & no \\
\bottomrule
\end{tabular}
\caption{The episode-sorting exhibit. The two capital-impairment onsets are exactly the two system-wide blow-outs; the two unimpaired onsets produced the smaller, short-end-concentrated moves. The constraint state at onset sorts the episode cross-section as the model predicts.}
\end{table}

\begin{figure}[h!]\centering
\includegraphics[width=0.92\textwidth]{fig11_basis_hkm.png}
\caption{Basis versus the HKM capital ratio; dashed: the estimated Spec-B threshold (bottom quintile). The two great widenings are the two excursions below the line.}
\end{figure}

\begin{figure}[h!]\centering
\includegraphics[width=0.78\textwidth]{fig12_constraint_test.png}
\caption{Calm versus stress loadings across specifications ($t$-statistics above bars).}
\end{figure}

\subsection{The forced-unwind surface}

The replay design, stated in full: entry on the first day of every month from 2004-08; the convergence trade is tracked on the \emph{daily} bond-level median with duration $D=8$ over a 24-month horizon; the position is stopped the first day the mark-to-market loss $D(\lambda_t-\lambda_0)$ exceeds the capital buffer $B$ (equivalently, the first crossing of the barrier $B/D$ above the entry spread); the surface reports the stop-out frequency by entry-year vintage and buffer. Daily monitoring is the empirical convention throughout, and Section~\ref{sec:struct} aligns the model to it.

\begin{table}[h!]
\centering\small
\begin{tabular}{lcccccccccc}
\toprule
Vintage & 2005 & 2007 & 2008 & 2011 & 2013 & 2016 & 2019 & 2021 & 2023 & 2025 \\
\midrule
Stop-out, 1\% & 100 & 100 & 100 & 67 & 92 & 8 & 100 & 83 & 0 & 0 \\
Stop-out, 2\% & 17 & 100 & 100 & 25 & 25 & 0 & 0 & 8 & 0 & 0 \\
\bottomrule
\end{tabular}
\caption{Stop-out frequency (\%) by entry vintage and capital buffer ($D=8$, 24-month horizon), bond-level median. The 2007--08 vintages are stopped out always, at both buffers; post-2013 the 2\% row is nearly empty --- the buffer is the binding economics, and the full surface is the structural model's calibration target.}
\end{table}

\subsection{The structural pass, corrected}\label{sec:struct}

Estimating \eqref{eq:ou} on the post-2013 median gives $\kappa=0.187$/month (half-life 3.7m), $\theta=16.6$\,bp, $\sigma=5.78$\,bp/m. The engine is rebuilt per the defect review: \emph{daily} barrier monitoring, aligned with the replay's convention; both horizons of Section~2.4 (the 24-month risk horizon with the MTM law for survivors, and bond maturity, which restores the note's mixture exactly); the stopped branch discounted at $Y_0=4\%$; and the certainty-equivalent condition computed exactly on the simulated terminal distribution.

\emph{The capital charge, derived (the note's $m\times\VaR$).} With $D\sigma=46$\,bp of monthly P\&L volatility (47 with the yield leg), the Basel ten-day 99\% VaR with the regulatory multiplier $m=3$ gives a charge of \textbf{2.26\%} of notional; the FRTB-style 97.5\% quantile gives \textbf{1.91\%}. The 2\% buffer that the replay identifies empirically is thus not a calibration: it is what the capital rule delivers under the estimated dynamics --- a free structural result that also makes the regulatory counterfactuals on $(m,\alpha)$ directly executable. Gaussian quantiles understate the charge given the daily-change kurtosis of 23.4; the fat-tail refinement raises it, in the direction favourable to the floor.

\begin{table}[h!]
\centering\small
\begin{tabular}{lcc}
\toprule
 & Model, daily monitoring & Empirical replay (daily) \\
\midrule
$\Pr(\text{stop}\le24\text{m})$, 1\% buffer & 68.8\% & 43.0\% \\
$\Pr(\text{stop}\le24\text{m})$, 2\% buffer & 8.8\% & 2.5\% \\
\bottomrule
\end{tabular}
\caption{Stop-out probabilities under the aligned (daily) convention; the 3-July table's monthly-monitoring values (58.9\%/6.5\%) are retained in the register for reference. The model over-predicts stop-outs at both buffers; the candidate explanations --- regime-dependent $\sigma$ (the post-2013 estimate averages calm and stress months), vintages entering below $\theta$, Gaussian versus fat-tailed shocks, and the omitted yield leg --- are competing, not adjudicated, and are exactly the targets of the SMM merge.}
\end{table}

\begin{table}[h!]
\centering\small
\begin{tabular}{llccccc}
\toprule
Horizon & Buffer & RN ($\mathbb{E}[X]\ge0$) & CE, RRA$=1$ & RRA$=2$ & RRA$=5$ & Hurdle 10\% \\
\midrule
$H=24$m & 1\% & 50.6 & 99.6 & $>$120 & $>$120 & 68.3 \\
$H=24$m & 2\% & 6.8 & 18.0 & 32.9 & $>$120 & 28.4 \\
$T=10$y & 1\% & 57.6 & 100.6 & $>$120 & $>$120 & 119.4 \\
$T=10$y & 2\% & 10.2 & \textbf{15.9} & \textbf{27.0} & $>$120 & 40.5 \\
\bottomrule
\end{tabular}
\caption{Model-implied floors $\lambda^*$ (bp/yr) under the note-faithful entry condition $\CE\ge W_0$ by risk aversion, with the risk-neutral participation floor and the 10\%-hurdle version as reference columns. Observed floor: 17.7 (median) -- 23.3 (10Y).}
\end{table}

The reading, in the review's own corrected terms. At the \emph{derived} $\approx$2\% charge and under the note's exact maturity mixture, the risk-neutral participation floor is 10\,bp and the CE floors at RRA $=1$ and $2$ are 15.9 and 27.0\,bp: \textbf{the observed floor of 17.7--23.3\,bp corresponds to relative risk aversion between one and two} --- squarely inside the standard range --- and the risk-aversion channel, not the expected-loss channel, is what generates most of it, exactly the note's one-sentence mechanism. The earlier ``brackets the observed floor'' phrasing is retired: on the aligned convention the observed median sits just below the RRA$=1$ value at the 24-month horizon and between the RRA$=1$ and RRA$=2$ values at maturity, and the sentence above is the precise statement. The 1\%-buffer rows, where risk aversion is first-order and the floors explode, are shown for completeness and are in any case excluded by the derived charge itself.

\begin{figure}[h!]\centering
\includegraphics[width=0.56\textwidth]{fig13_structural_pass.png}
\caption{Stop-out probabilities, model versus data --- \emph{monthly-monitoring version of 3 July}, retained for comparability; the aligned daily figures are in the table above.}
\end{figure}

\subsection{Replicating FLL Figure 2, and the per-pair diagnostic}\label{sec:fllrep}

The paper's construction is finally confronted with FLL's own published object: their Figure~2, the daily cross-sectional average mispricing in basis points, weighted by the notional of each TIPS issue, over July 23, 2004 -- November 19, 2009. The replication uses the 107-CUSIP panel with per-issue Amounts Issued (matched 107/107) as weights. Result: window mean \textbf{40.4\,bp weighted} (39.8 equal-weighted) against their 54.5; peak \textbf{153.5\,bp on 2008-12-31} against their $\sim$175 at the Lehman apex; time profile identical (Figure~\ref{fig:fllext}). Two findings follow. First, \emph{weighting is second-order} (0.6\,bp) --- consistent with FLL's own Figure~3 showing no maturity pattern in mispricing --- so the residual gap sits in the pair construction upstream. Second, a correction is recorded: the deflation floor is \emph{not} the explanation --- FLL explicitly abstract from the floor (their Section~D) and note adjusting for it would make mispricing \emph{larger}, so a panel that also ignores it inherits no gap from that source.

The gap is then localised with the per-pair diagnostic: FLL's Table~III (29 pairs with per-pair mean, SD, extremes and $N$; extracted to \texttt{fll\_pairs\_tableIII.csv}) is matched to panel columns by maturity, 24/29 matched. The verdict is sharp:

\begin{table}[h!]
\centering\small
\begin{tabular}{lccc}
\toprule
Pair group (maturity mismatch) & Pairs matched & Mean diff, panel $-$ FLL (bp) & Reading \\
\midrule
31-day (Feb-cycle, 2015--2029) & 11 & $\mathbf{+2.8}$ & \textbf{certified}: panel $\approx$ FLL within $\pm$3--4\,bp \\
15-day (2007--2014, Jan/Mar cycles) & 9 & $-23.1$ & short-end construction gap \\
0-day (exact matches, 2009--2012) & 4 & $-29.2$ & short-end construction gap \\
\bottomrule
\end{tabular}
\caption{Per-pair diagnostic (full table: Appendix, \texttt{fll\_pair\_diagnostic.csv}). Corr(diff, FLL mean) $=-0.45$: the panel undershoots most where FLL measure most. The smoking gun is in the observation counts: for several short pairs the panel has $N\!\approx\!1{,}327$ days where FLL's pair has $N=91$--$395$ --- FLL's pair exists only from the issuance of the \emph{specific matched Treasury} (e.g., the 1.125\% Jan-2012, issued 2009), so the panel is pairing those TIPS against a \emph{different} nominal over the earlier period. Combined with the absence of FLL's seasonal adjustment on fractional-maturity swap legs (their Appendix A1; our measured 52\,bp two-year amplitude), this fully accounts for the short-end gap. Consequence: the long end of the construction is certified at the individual-pair level; a dedicated FLL engine for the $\sim$15 short pairs (exact Table-III Treasuries, per-cash-flow seasonally adjusted swaps, STRIPS coupon matching, synthetic-yield-to-matched-maturity) is a designated task, with Table~III itself as the 29-row validation suite.}
\end{table}

\begin{figure}[h!]\centering
\includegraphics[width=0.94\textwidth]{fll_fig2_extended.png}
\caption{FLL Figure 2, replicated (notional-weighted, their style) and extended to 2026; dashed line: end of their sample (Nov-2009). Everything to the right is the out-of-sample of the original paper --- and the subject of this one.}\label{fig:fllext}
\end{figure}

\section{Established versus Open}

\emph{Established:} the floor's existence, level, and dynamics on a triple-validated construction, now certified pair-by-pair at the long end against FLL's Table~III; the staircase with dated breaks; the certainty-equivalent floor at the derived capital charge; the seasonal decomposition of the short end; the event phenomenology with lags; the dispersion signature of March 2020; the state-dependent constraint interaction ($t=4.1$) with the episode sorting; the survival surface; and a first structural pass in which the model, with data-estimated dynamics and economically pinned frictions, lands on the observed floor. \emph{Open, stated plainly:} the formal threshold-versus-linear separation remains short of conventional significance on U.S. data alone (sup-$F$ $p\approx0.25$; pre-specified Chow $p=0.10$/$0.15$) and the leave-one-episode-out check shows the interaction is substantially GFC-driven; the 2024 short-end closure is \emph{provisional} (twenty post-break months; deseasonalised mean 2.8\,bp rather than the raw 0.3); neither candidate international extension is free of qualification --- the euro episode (2011--12) is not a clean third impairment until sovereign and redenomination risk are explicitly controlled (a Bund\texteuro i-versus-OAT\texteuro i differential with quanto-CDS controls), while the UK LDI episode (October 2022) is clean of the sovereign confound but is primarily a forced unwind of \emph{end investors} under collateral calls rather than an impairment of dealer capital, so it bears on the barrier mechanism and H6 more directly than on the H3 capital channel; the short-maturity pairs and March 2020 at benchmark maturities require the dedicated FLL engine; any negative-floor/2021 claim remains gated by the 52\,bp seasonal and by unavailable specialness data; the transaction-cost band is assumed (zero), conservative for existence but leaving the no-trade-band interpretation of 20\,bp unresolved; the full SMM (June engine, GKR covariance, CE and correlated-yield refinements, dispersion moment) and the policy counterfactuals are designed but not yet run in the unified pipeline.

\section{Road to Completion and Publication Strategy}

The remaining programme, in execution order: (i) merge the June \texttt{tips\_treasury\_arb} SMM engine with the July moment set --- the survival surface plus the Spec-B loading pair, GKR block-bootstrap moment covariance, CE and correlated-yield refinements, the dispersion moment from \eqref{eq:mixture} (no new data); (ii) the dedicated FLL engine for the short pairs and March 2020 (ISIN-level prices for the Table-III bonds, index ratios, bid--ask), with the indexation-schedule seasonal; (iii) the regulatory-calendar test of the 2010/2017 staircase and the policy counterfactuals on $(m,\alpha,F)$; (iv) the international linker cross-section --- the power play for the threshold test and the single addition with the highest expected effect on the paper's tier --- built from two markets supplying the impairment episodes the U.S. sample lacks: the \emph{euro} bases (2011--12), designed from the outset with the sovereign-risk control the review demands (the Bund\texteuro i-versus-OAT\texteuro i \emph{differential} basis, quanto-CDS and redenomination controls, the pooled claim made only conditional on them), and the \emph{UK} index-linked gilt basis (the October 2022 LDI dislocation), which is orthogonal to both U.S. clusters, carries no sovereign confound, and --- index-linked gilts having no deflation floor on principal --- doubles as a control for the floor-option wedge that TIPS and euro linkers share; scope is capped at these two, Japanese linkers excluded as thin and yield-curve-control distorted; (v) robustness per the original roadmap (dual construction switches, specification curves, placebos on macro fundamentals, costs stressed). Targets: \emph{Review of Finance}/\emph{JFQA} if (i)+(iv) land, with \emph{JEF}/\emph{JBF} as the secured field tier on material in hand; a practitioner distillation for the \emph{Journal of Fixed Income} only after the main submission.

\appendix

\section{Reconciliation with the June 2026 Detailed Project Description}

The 15-page \emph{Detailed Project Description and Preliminary Results} (17--19 June 2026), produced with the \texttt{tips\_treasury\_arb/} codebase, is this document's direct ancestor. The reconciliation, item by item:

\begin{table}[h!]
\centering\small
\setlength{\tabcolsep}{3.5pt}\begin{tabular}{p{5.3cm}p{3.1cm}p{5.5cm}}
\toprule
June result & Status & July verdict \\
\midrule
Floor: post-2013 median 17.7\,bp (pre 37.6; peak 153.4, Nov-2008) & \textbf{Confirmed} & Identical on exact legs and bond level \\
Break at 2010, not 2013 (2010-09, sup-$F$ 2533, median constr.) & \textbf{Confirmed, refined} & 2010-12 (sup-$F$ 160, monthly exact legs) + \emph{second} break 2017-12 + 2024-10 short-end closure: a staircase \\
OU: $\theta=17.2$, half-life 1.9\,m & Confirmed / \emph{corrected} & $\theta=16.6$; HL 3.7--4.6\,m on the longer exact samples (window convention) \\
HKM: insignificant unconditionally, $-26.2$ ($p<0.05$) in top-quintile-vol months & \textbf{Confirmed, redesigned} & Spec A: $-24.5$ [$-2.6$] in the HKM-low state --- near-identical magnitude from an independent design; plus Spec B (7.15 [4.1]) and the episode-sorting table \\
Bidirectionality: 76 negative days 2021--22, min $-7.5$ (27-05-2021), symmetric-model RQ & \emph{Downgraded / parked} & Largely the 52\,bp two-year seasonal; residual unidentifiable without security-level specialness; signed model held in reserve \\
March 2020 only $\sim$2$\sigma$ at the median (identification caveat) & \emph{Resolved as a discovery} & A \emph{dispersion} event: IQR peak 2020-03-16, 90th pct 73\,bp --- prediction \eqref{eq:mixture} \\
Episode detection: only GFC and June-2013 taper clear 3$\sigma$ & Superseded & Replaced by the episode-HKM sorting (a constraint-based, not amplitude-based, classification) \\
Daily-change kurtosis 23.4 (jump risk present) & Retained & Feeds the fat-tail robustness of the first-passage engine \\
Full SMM (McFadden/Duffie--Singleton) in \texttt{tips\_treasury\_arb} & \textbf{To merge} & Run with the July moments (surface + loading pair) and GKR covariance --- next step (i) \\
Policy counterfactuals (RQ5) & Pending & Scheduled with step (iii) \\
GIRF (Pesaran--Shin) and quantile (Koenker--Bassett) batteries & Robustness queue & Partially superseded by the threshold design; retained as robustness \\
\bottomrule
\end{tabular}
\end{table}

\section{Register of Every Test Executed to Date}

All tests on the TIPS project, in execution order, with script and headline output. (The regime-consistent-resampling pilots are archived separately in \emph{Pilot Report v4}.)

\begin{table}[h!]
\centering\footnotesize
\setlength{\tabcolsep}{3.5pt}\begin{tabular}{p{0.8cm}p{5.4cm}p{3.2cm}p{4.5cm}}
\toprule
ID & Test & Script & Headline output \\
\midrule
T1--T2 & Subperiod means, NW $t$; term structure (proxy legs) & \texttt{tips\_pilot.py} & Post-2013 floor 12.7--23.8\,bp, $t$ 7.6--13.2; proxy inversion (later resolved) \\
T3 & Single-break SSR grid, per maturity & \texttt{tips\_pilot.py} & First-pass 2019--21 ``wave'' (later exposed as truncation artefact) \\
T4 & ADF; OU half-life, bootstrap CI & \texttt{tips\_pilot.py} & 10/20Y stationary, HL 1.9\,m (2020s window) \\
T5--T6 & Stress correlations, events, resting-vs-MOVE & \texttt{tips\_pilot.py} & Event 5--7$\times$ moves; lag discovered; linear $\approx$ 0 \\
--- & Coverage diagnostics (series start dates) & inline & USGG20YR from 2020-05; 7YR from 2009 $\Rightarrow$ four-ticker fix \\
D1--D3 & Exact-leg means; FLL-window benchmark & \texttt{tips\_pilot2.py} & 46.8 vs 54.5\,bp; floor 16--24, $t$ to 24.2 \\
D2 & Breakeven cross-check, 5 maturities & \texttt{tips\_pilot2.py} & 10Y: $-0.55$\,bp, corr 0.996 \\
D5 & Break adjudication, 1st + 2nd breaks & \texttt{tips\_pilot2.py} & 2010-12 (supF 160); 2017-12; 2024-10 closure \\
D6 & ADF/OU + month-dummy seasonality & \texttt{tips\_pilot2.py} & 52.2\,bp amplitude at 2Y; ADF flip 0.65$\to$0.000 \\
D7 & Events on exact legs & \texttt{tips\_pilot2.py} & GFC 144--154; SVB peaks Jul--Aug (lag) \\
D8--D9 & Lead-lag correlogram; local projections; rolling floor & \texttt{tips\_pilot2.py} & All linear $|t|\le1.2$; $\beta=-0.014$ [$-0.4$] \\
--- & Mar-2020/GFC addendum, 2Y (exact legs) & inline & 2Y $-1.3\to68.2$ (peak 19-03), half-rev 8 days; GFC 2Y 321.9 \\
B1--B2 & CUSIP panel coverage; validation vs curve & \texttt{cusip\_analysis.py} & corr 0.926 level, 0.591 $\Delta$monthly \\
B3--B4 & Mar-2020 bond level; maturity split & \texttt{cusip\_analysis.py} & IQR 5.3$\to$18.4 (peak 16-03); 90th pct 73.3 \\
B5 & Threshold, contemporaneous, market proxies & \texttt{cusip\_analysis.py} & All $p>0.14$, perverse signs: retired \\
B6 & Forced-unwind replay, vintage $\times$ buffer & \texttt{cusip\_analysis.py} & 2007--08: 100\%/100\%; post-13 at 2\%: $\approx$0 \\
Add.1 & Threshold, lag-aware, market proxies & inline & CDS perverse ($p=0.023$); LOIS/MOVE $t\approx3$, $p=0.21/0.36$ \\
Add.2--3 & Dispersion by episode; $\Delta$IQR loadings & inline & Event-concentrated; $|t|\le0.8$ \\
A--D & Constraint test, frozen design (HKM/PD/TFF) & \texttt{constraint\_test.py} & Spec B 7.15 [4.1] vs 0.63; C/D uninformative \\
--- & Block-wild refinement; episode-HKM table & inline & $p$ stable 0.24--0.26; 3\%/8\% vs 32\%/37\% \\
--- & OFR quarterly interaction (repo $\times$ RV-NAV) & \texttt{constraint\_test.py} & Uninformative, $n=52$ \\
S1 & Structural pass: OU $\to$ first-passage $\to$ floor & \texttt{structural\_pass} & Stop-out 58.9/6.5 vs 43/2.5\%; $\lambda^*$ 18--28 vs 17.7--23.3 \\
F1 & FLL Fig.\,2 replication, notional-weighted & \texttt{plot\_fll\_figure2.py} & 40.4 vs 54.5; peak 153.5 weighted / 151.3 equal (31-12-2008) \\
S2 & Engine v2: daily monitoring, two horizons, CE & \texttt{tier1\_fixes} & stop 68.8/8.8\%; floors: obs.\ between RRA 1--2 at 2\% \\
S3 & Capital charge derived ($m\times\VaR$) & \texttt{tier1\_fixes} & 2.26\% (Basel) / 1.91\% (FRTB) \\
S4 & Chow at pre-specified quintile; LOO; composition & \texttt{tier1\_fixes} & $p=0.097$/0.145; ex-GFC collapses (2.46 [1.8]) \\
S5 & Dispersion: uniform pre rule; within-bucket IQR & \texttt{tier1\_fixes} & 2022 pre 16.2; $<$5y IQR 5.3$\to$43.8 vs $\ge$5y 1.3$\to$4.4 \\
S6 & 2Y post-2024 deseasonalised mean & \texttt{tier1\_fixes} & 2.8\,bp ($n=20$m): closure \emph{provisional} \\
F2 & Table-III per-pair diagnostic & inline & Long pairs $+2.8$; short $-23/-29$; $N$ smoking gun \\
\bottomrule
\end{tabular}
\end{table}

\section{Imported Evidence and Reinstated Caveats}

Three evidence blocks asserted in the main text are \emph{imported} from the 3-July Comprehensive Report and are labelled as such pending re-verification in the merged pipeline: the 19-variable $\times$ 3-frequency spurious-trend battery (the systematic dissolution of levels relationships in differences that underwrites the ``shocks, not constraints'' reading of \S\ref{sec:constraint}); the composite-factor ex-GFC collapse across the full term structure ($R^2$ $0.42\to0.02$ at ten years, analogues at all maturities), which generalises the linear-emptiness result; and the June SMM/GKR robust-inference block ($J$ $7.56\to1.21$ under block-bootstrap moment covariance), citable but unverified until the merge. Three caveats from earlier versions are reinstated verbatim in spirit: the 2024 short-end result rests on twenty post-break months and is provisional; the levels-versus-changes discipline cuts \emph{both} ways (it dissolves spurious support and spurious refutation alike); and the accumulating battery of tests warrants a multiple-testing acknowledgement, with the pre-specified objects (the frozen design; the bottom quintile) carrying the inferential weight precisely because they were fixed ex ante.

\begin{thebibliography}{99}
\bibitem[Alili et al.(2005)]{alili2005} Alili, L., Patie, P., Pedersen, J.L., 2005. Representations of the first hitting time density of an Ornstein--Uhlenbeck process. \emph{Stochastic Models} 21, 967--980.
\bibitem[Christensen and Gillan(2022)]{cg2022} Christensen, J.H.E., Gillan, J.M., 2022. Does quantitative easing affect market liquidity? \emph{Journal of Banking and Finance} 134, 106349.
\bibitem[d'Amico et al.(2018)]{dkw2018} d'Amico, S., Kim, D.H., Wei, M., 2018. Tips from TIPS. \emph{Journal of Financial and Quantitative Analysis} 53, 395--436.
\bibitem[Duffie(2010)]{duffie2010} Duffie, D., 2010. Asset price dynamics with slow-moving capital. \emph{Journal of Finance} 65, 1237--1267.
\bibitem[Fleckenstein and Longstaff(2024)]{fl2024} Fleckenstein, M., Longstaff, F.A., 2024. Treasury richness. \emph{Journal of Finance} 79, 2797--2844.
\bibitem[Fleckenstein et al.(2014)]{fll2014} Fleckenstein, M., Longstaff, F.A., Lustig, H., 2014. The TIPS--Treasury bond puzzle. \emph{Journal of Finance} 69, 2151--2197.
\bibitem[G\^arleanu and Pedersen(2011)]{gp2011} G\^arleanu, N., Pedersen, L.H., 2011. Margin-based asset pricing and deviations from the law of one price. \emph{Review of Financial Studies} 24, 1980--2022.
\bibitem[Gospodinov et al.(2014)]{gkr2014} Gospodinov, N., Kan, R., Robotti, C., 2014. Misspecification-robust inference in linear asset-pricing models. \emph{Review of Financial Studies} 27, 2139--2170.
\bibitem[Gromb and Vayanos(2010)]{gromb2010} Gromb, D., Vayanos, D., 2010. Limits of arbitrage. \emph{Annual Review of Financial Economics} 2, 251--275.
\bibitem[He et al.(2017)]{hkm2017} He, Z., Kelly, B., Manela, A., 2017. Intermediary asset pricing. \emph{Journal of Financial Economics} 126, 1--35.
\bibitem[He et al.(2022)]{hns2022} He, Z., Nagel, S., Song, Z., 2022. Treasury inconvenience yields during the COVID-19 crisis. \emph{Journal of Financial Economics} 143, 57--79.
\bibitem[Kondor(2009)]{kondor2009} Kondor, P., 2009. Risk in dynamic arbitrage. \emph{Journal of Finance} 64, 631--655.
\bibitem[Linetsky(2004)]{linetsky2004} Linetsky, V., 2004. Computing hitting time densities for CIR and OU diffusions. \emph{Journal of Computational Finance} 7, 1--22.
\bibitem[Liu and Longstaff(2004)]{liulongstaff2004} Liu, J., Longstaff, F.A., 2004. Losing money on arbitrage. \emph{Review of Financial Studies} 17, 611--641.
\bibitem[Mitchell and Pulvino(2012)]{mp2012} Mitchell, M., Pulvino, T., 2012. Arbitrage crashes and the speed of capital. \emph{Journal of Financial Economics} 104, 469--490.
\bibitem[Rebonato(2026)]{rebonato2026} Rebonato, R., 2026. The Idea. Unpublished note, EDHEC Business School.
\bibitem[Shleifer and Vishny(1997)]{shleifervishny1997} Shleifer, A., Vishny, R.W., 1997. The limits of arbitrage. \emph{Journal of Finance} 52, 35--55.
\bibitem[Tuckman and Vila(1992)]{tuckmanvila1992} Tuckman, B., Vila, J.-L., 1992. Arbitrage with holding costs. \emph{Journal of Finance} 47, 1283--1302.
\end{thebibliography}

\end{document}
